\section{Methods}
\subsection*{Model description}

Here we employ the NEMO numerical platform (release 4.0.5; \citealp{NEMO_man}) to build a 1/60$^\circ$ pan-Arctic ocean-sea ice configuration named SEDNA. The ocean core (NEMO-OPA) solves the primitive equations using finite differences on an Arakawa C-grid mesh, and it is coupled to the SI3 sea-ice component (NEMO-SI3; \citealp{SI3_man}). The pan-Arctic domain covers the Bering Strait on the Pacific side and extends southward to 56$^\circ$N in the Subpolar North Atlantic and to the Baltic Sea entrance. The ocean model incorporates a linear free surface, utilizes a 3rd order flux form scheme for momentum advection, and employs a 4th order Flux Corrected Transport (FCT) for tracer advection. Additionally, the sea-ice representation uses an Elasto-Visco-Plastic (EVP) rheology with a 5-categories sea ice thickness distribution and a landfast ice parameterization to better simulate sea-ice behavior above shelves.

The NCAR bulk formula from \citet{Large_global_2009} is used to calculate air-sea fluxes based on the hourly ERA5 atmospheric data for near-surface variables \citep{Hersbach_ERA5_2020}. The three lateral boundaries are constrained with daily GLORYS12V1 Re-analysis data \citep{GLORYS12}. Monthly freshwater fluxes from river discharges are combined with Greenland land ice melt data \citep{HuEA_rivers_2019}. The simulation extends over seven years (2009 to 2015) and starts from rest, utilizing initial conditions based on the World Ocean Atlas 2009 temperature/salinity and mean January 2009 ice state from PIOMAS re-analysis \citep{Locarnini_data_2010, PIOMAS}. A weak restoring toward the World Ocean Atlas climatological sea surface salinity is applied, but only over the ice-free regions. Preliminary assessment against available observations and re-analyses data shows that the simulation reasonably captures key ocean and sea-ice characteristics.

\subsection*{Identification of leads}

The leads from SEDNA are identified by using the algorithm proposed by \citet{Hutter_identification_2019}. This method uses the total deformation rate defined as: 
\begin{equation}
    T_d = \left[\left(\underbrace{\frac{\partial u}{\partial x} +\frac{\partial v}{\partial y}}_{Divergence}\right)^2 + \underbrace{\left(\frac{\partial u}{\partial x} - \frac{\partial v}{\partial y}\right)^2+ \left(\frac{\partial u}{\partial y} + \frac{\partial v}{\partial x}\right)^2}_{Shear}\right]^{1/2},
\end{equation}
where $u$ and $v$ are the sea ice velocities. The deformation rate varies spatially depending on the background deformation and the sea ice properties \citep{Reiser_LKF_2020}, where the largest deformation rates occur along the boundaries of ice floes where the leads are found. Therefore, to identify the leads, the local maxima is identified using a difference of Gaussian filter (DoG filter) of the deformation rate field. After this, a binary map is created with the mask of identified lead. Note that after this step, \citet{Hutter_identification_2019} reduce the leads width to a line to then track the leads over time. Since our study focuses on the ocean-atmosphere interactions, we only use the identification algorithm to produce a mask of leads during 2014 for the Arctic, thus we retain the original lead width computed from the deformation rate. Finally, the lead mask is then used to quantify and contrast the impact of leads compared to the ice covered Arctic, the ice-pack (ice concentration $> 80\%$), and the MIZ (0 - 80\% ice concentration).

\subsection*{Buoyancy flux and water mass transformation}
 
Heat fluxes (shortwave radiation, longwave radiation, latent heat, and sensible heat) and freshwater fluxes (evaporation, precipitation, runoff, and sea ice melt/growth) at the ocean surface result in a buoyancy forcing capable of changing the density of the ocean surface. Here, the buoyancy flux is computed as:

\begin{equation}
    B_f = \underbrace{\frac{g \alpha Q_{net}}{\rho_0 C_p}}_{Net\ surface\ heat\ flux} - \underbrace{\frac{g \beta F_{net} SSS }{\rho_0}}_{Net\ surface\ freshwater\ flux},
\end{equation}

where $Q_{net}$ the net heat flux at the sea surface, $\alpha$ the thermal expansion coefficient, $SSS$ the sea surface salinity, $F_{net}$ the net freshwater flux that includes evaporation, precipitation, runoff, and freezing minus melting flux, $\beta$ the haline contraction coefficient, $g = 9.81 m/s^2$, and the ocean density $\rho_0 = 1026 kg/m^3$. 
Note that using this sign convention, positive values in the buoyancy forcing decrease the surface density making the ocean surface more buoyant (stable) and negative values increase the surface density. The area-weighted buoyancy ($ \left<\right>$) fluxes are computed following:

\begin{equation}
    \left<B_f (t)\right>  = \frac{\sum_x \sum_y \left( B_f(t,x,y) * Area(x,y) * Mask_{_R}(t,x,y)\right)}{\sum_x \sum_y * Area(x,y) * Mask(t,x,y)},
\end{equation}
where the $Area$ is the area of the grid cells, and the $Mask_{_R}$ can be the mask of MIZ, the Pack, the MIZ leads, the ice pack leads, and the ice cover that varies in space and time. Note that the divisor $Mask$ corresponds to the full ice covered area (including leads).

\citet{Nurser_WMT_1999} and  \citet{Walin_WMT_1982} proposed the water mass transformation framework that quantifies the rate at which a given water mass is made lighter (negative transformation rate) or denser (positive) due to surface fluxes. The surface water mass transformation is computed as:

\begin{equation}
    \Omega(\sigma_k) = - \frac{1}{\sigma_{k+1}-\sigma_k}\iint_AB_fdA,
\end{equation}

where $\sigma$ is potential density and the indices $k$ represents the density bin number. The computation was performed with 172 density bins of $0.1 kg/m^3$ within the density range of $15.15 kg/m^3$ - $32.25 kg/m^3$. The transformation rate $\Omega$ is spatially decomposed into different sea-ice masks such as the ice covered, the pack ice, the MIZ, and the leads in each of these regions. 